% Slides. Overlays of one frame keep the same picture. Run lualatex twice.
\documentclass[aspectratio=169]{beamer}
\usepackage{limen}
\begin{document}
\begin{frame}[plain]
  \limencover{How light earns its colour}{A short talk}{Seminar, November 2026}
\end{frame}
\begin{frame}{Below the threshold}
  \begin{itemize}
    \item<1-> The beam is thin and almost white
    \item<2-> The field is dark and the lines are only beginning to show
    \item<3-> Overlays of one frame keep the same picture
  \end{itemize}
\end{frame}
\begin{frame}{The curve}
  Output stays low until the threshold, then rises almost in a straight line.
\end{frame}
\begin{frame}{The threshold}
  \limenthreshold
  You can put the threshold on the frame where your argument turns.
\end{frame}
\begin{frame}{Above it}
  \begin{enumerate}
    \item The spindles grow wider
    \item The field gains colour
  \end{enumerate}
\end{frame}
\begin{frame}{Almost there}
  A strip of spectrum: \limenspectrum[4cm]
\end{frame}
\begin{frame}{The end}
  The last frame closes with a small fan that leaves through the edge of the sheet.
\end{frame}
\end{document}
